\documentclass[preview]{standalone}

\usepackage{amsmath}
\usepackage{amssymb}
\usepackage{stellar}
\usepackage{enumitem}
\usepackage{definitions}

\begin{document}

\id{linearalgebra-exercises-batch-13}
\genpage

\section{Exercises}

\begin{snippetexercise}{linear-algebra-batch-13-ex-1}{}
    Given \(c_1, \dots, c_n \in \mathbb{F}\), \(c_i > 0\) for every \(i = 1, \dots, n\),
    we define a map
    \[\langle \cdot, \cdot \rangle \colon \mathbb{F}^n \times \mathbb{F}^n \fromto \mathbb{F}, \quad \langle z, w \rangle \triangleq c_1 z_1 \overline{w_1} + \dots + c_n z_n \overline{w_n}.\]
    Verify that \((\mathbb{F}^n, \langle \cdot, \cdot \rangle)\) is an inner product space.
\end{snippetexercise}

\begin{snippetsolution}{linear-algebra-batch-13-ex-1-sol}{}
    \todo
\end{snippetsolution}

\begin{snippetexercise}{linear-algebra-batch-13-ex-2}{}
    Let \(V\) be an \(n\)-dimensional \(\mathbb{F}\)-vector space and let
    \(S \in \matrices_{n \times n}(\mathbb{F})\) be a symmetric
    (\(\mathbb{F} = \realnumbers\)) or Hermitian 
    (\(\mathbb{F} = \complexnumbers\)) and positive definite \matrix.
    Verify that the \function
    \[\langle \cdot, \cdot \rangle \colon V \times V \fromto \mathbb{F}, \quad \langle v, w \rangle \triangleq \Phi_{\mathcal{B}}(v)^t S \overline{\Phi_{\mathcal{B}}(w)}\]
    defines an inner product on \(V\)
    (the vectors \(\Phi_{\mathcal{B}}(v)\) are interpreted as column vectors).
\end{snippetexercise}

\begin{snippetsolution}{linear-algebra-batch-13-ex-2-sol}{}
    \todo
\end{snippetsolution}

\begin{snippetexercise}{linear-algebra-batch-13-ex-3}{}
    Let \(n \geq 2\).
    On \(\mathcal{P}_{n-1}(\realnumbers)\) consider the map
    \(\langle \cdot, \cdot \rangle \colon \mathcal{P}_{n-1}(\realnumbers) \times \mathcal{P}_{n-1}(\realnumbers) \fromto \realnumbers\), \(\langle p, q \rangle \triangleq \int_0^1 pq\).
    \begin{enumerate}[label=\Alph*]
        \item Prove that this is an inner product.
        \item Prove that the matrix associated with \(\langle \cdot, \cdot \rangle\) with respect to the basis \(\{1, x, \dots, x^{n-1}\}\) is the \emph{Hilbert matrix}, \(H_n\), whose \((i, j)\)-th entry is \(\frac{1}{i+j-1}\).
        \item Conclude that the Hilbert matrix is positive definite.
    \end{enumerate}
\end{snippetexercise}

\begin{snippetsolution}{linear-algebra-batch-13-ex-3-sol}{}
    \todo
\end{snippetsolution}

\begin{snippetexercise}{linear-algebra-batch-13-ex-4}{}
    Show that the following pairs \emph{are not} inner product spaces:
    \begin{enumerate}[label=\Alph*]
        \item \(V = \realnumbers^2\), \(\langle (x_1, x_2), (y_1, y_2) \rangle \triangleq |x_1 y_1| + |x_2 y_2|\).
        \item \(V = \realnumbers^2\), \(\langle (x_1, x_2), (y_1, y_2) \rangle \triangleq 2x_1 y_2 + 3x_2 y_1\).
        \item \(V = \mathcal{P}_3(\realnumbers)\), \(\langle p, q \rangle \triangleq p(0)q(0) + p(1)q(1) + p(-1)q(-1)\).
    \end{enumerate}
\end{snippetexercise}

\begin{snippetsolution}{linear-algebra-batch-13-ex-4-sol}{}
    \todo
\end{snippetsolution}

\begin{snippetexercise}{linear-algebra-batch-13-ex-5}{}
    In the real vector space
    \[V = \left\{ \begin{pmatrix} a & b & 0 \\ c & a & b \\ 0 & c & a \end{pmatrix} \suchthat a, b, c \in \realnumbers \right\} \subset \matrices_{3 \times 3}(\realnumbers)\]
    consider the inner product
    \(\langle A, B \rangle \triangleq \trace(AB^t)\).
    Find an orthonormal basis of \((V, \langle \cdot, \cdot \rangle)\).
\end{snippetexercise}

\begin{snippetsolution}{linear-algebra-batch-13-ex-5-sol}{}
    \todo
\end{snippetsolution}

\begin{snippetexercise}{linear-algebra-batch-13-ex-6}{}
    In \(\matrices_{2 \times 2}(\complexnumbers)\) consider the subset \(V\) of traceless Hermitian matrices: \[V = \{A \in \matrices_{2 \times 2}(\complexnumbers) \suchthat A = \bar{A}^t, \trace(A) = 0\}.\]
    \begin{enumerate}[label=\Alph*]
        \item Verify that \(V\) is a real vector subspace of \(\matrices_{2 \times 2}(\complexnumbers)\).
        \item Find a basis of \(V\), orthonormal with respect to the restriction to \(V\) of the inner product \(\langle \cdot, \cdot \rangle \colon \matrices_{2 \times 2}(\complexnumbers) \times \matrices_{2 \times 2}(\complexnumbers) \fromto \complexnumbers\), \(\langle A, B \rangle \triangleq \trace(A\bar{B}^t)\).
    \end{enumerate}
\end{snippetexercise}

\begin{snippetsolution}{linear-algebra-batch-13-ex-6-sol}{}
    \todo
\end{snippetsolution}

\begin{snippetexercise}{linear-algebra-batch-13-ex-7}{}
    Consider the bilinear map \(\langle \cdot, \cdot \rangle \colon \realnumbers^3 \times \realnumbers^3 \fromto \realnumbers\) defined by \[\langle (x_1, x_2, x_3), (y_1, y_2, y_3) \rangle \triangleq x_1y_1 + 2x_2y_2 - x_2y_3 - x_3y_2 + 2x_3y_3.\]
    \begin{enumerate}[label=\Alph*]
        \item Verify that this is an inner product on \(\realnumbers^3\).
        \item Find an orthonormal basis of \((\realnumbers^3, \langle \cdot, \cdot \rangle)\).
    \end{enumerate}
\end{snippetexercise}

\begin{snippetsolution}{linear-algebra-batch-13-ex-7-sol}{}
    \todo
\end{snippetsolution}

\begin{snippetexercise}{linear-algebra-batch-13-ex-8}{}
    Consider the inner product space \((\mathcal{P}_2(\realnumbers), \langle \cdot, \cdot \rangle)\), with \(\langle p, q \rangle \triangleq \int_0^1 pq\).\\
    \begin{enumerate}[label=\Alph*]
        \item Find an orthonormal basis of \((\mathcal{P}_2(\realnumbers), \langle \cdot, \cdot \rangle)\).
        \item Find a polynomial \(q \in \mathcal{P}_2(\realnumbers)\) such that \[p\left(\frac{1}{2}\right) = \int_0^1 pq \quad \forall p \in \mathcal{P}_2(\realnumbers).\]
    \end{enumerate}
\end{snippetexercise}

\begin{snippetsolution}{linear-algebra-batch-13-ex-8-sol}{}
    \todo
\end{snippetsolution}

\end{document}